\begin{table}[htbp]
  \centering
  \caption{Polynomial approximation to the firm-size wage profile}
  \label{tab:polynomial-size-profile}
  \small
  \begin{tabular}{lccc}
    \toprule
    Sample & $\ln(\widetilde{Size})$ & $[\ln(\widetilde{Size})]^2$ & Implied 101+/solo premium \\
    \midrule
    Raw & 0.241*** & -0.008 & 175.8\%*** \\
     & (0.026) & (0.007) &  \\
    Full sample & 0.190*** & -0.024*** & 41.0\%*** \\
     & (0.035) & (0.006) &  \\
    Formal workers & -0.038 & 0.011 & 8.6\% \\
     & (0.039) & (0.006) &  \\
    Informal workers & 0.260*** & -0.037*** & 46.9\%*** \\
     & (0.026) & (0.005) &  \\
    \bottomrule
  \end{tabular}
  \vspace{0.3em}
  \begin{minipage}{0.95\textwidth}
  \footnotesize
  Notes: The table reports quadratic specifications in assigned log firm size. The implied 101+/solo premium evaluates the fitted polynomial at 150 workers relative to solo work, for which $\ln(\widetilde{Size})=0$. Standard errors for the polynomial coefficients are clustered by sector and reported in parentheses. Stars in the premium column are based on the underlying log 101+/solo contrast. Because GEIH reports firm size in bins, the estimates are an approximation to the continuous size-wage profile and should be read as a robustness and shape check rather than as a replacement for the categorical estimates. Significance levels: * $p<0.10$, ** $p<0.05$, *** $p<0.01$.
  \end{minipage}
\end{table}
